% !TEX TS-program = pdflatex
% !TEX encoding = UTF-8 Unicode
% !BIB TS-program = bibtex
%
% Cycling Poverty in France: A Triple-Penalty Diagnostic on 34,858 Communes.
% Authors : R. Foss\'e \& G. Pallares (CESI LINEACT, 2025--2026).
%
\documentclass[11pt,a4paper]{article}

\usepackage[T1]{fontenc}
\usepackage[utf8]{inputenc}
\usepackage{lmodern}
\usepackage[margin=2.5cm]{geometry}
\usepackage{microtype}
\usepackage{csquotes}

\usepackage{amssymb,amsmath,amstext}

\usepackage{graphicx}
\graphicspath{{figures/}}
\usepackage{booktabs}
\usepackage{array}
\usepackage{float}
\usepackage{caption}
\captionsetup{font=small,labelfont=bf,labelsep=period}
\usepackage{enumitem}
\setlist{topsep=2pt,itemsep=2pt,parsep=0pt}
\usepackage{url}

\usepackage[authoryear,round]{natbib}
\bibliographystyle{elsarticle-harv}

\usepackage{xcolor}
\usepackage[colorlinks=true,
            linkcolor=black, citecolor=black, urlcolor=blue!50!black,
            breaklinks=true]{hyperref}

\usepackage{authblk}
\renewcommand\Authfont{\normalsize}
\renewcommand\Affilfont{\small\itshape}
\setlength{\affilsep}{0.6em}

\newcommand{\IMD}{\mathrm{IMD}}
\newcommand{\IES}{\mathrm{IES}}
\newcommand{\TPI}{\mathrm{TPI}}

% =========================================================================
\title{\Large\bfseries Cycling poverty in France:\\
\large a triple-penalty mobility-justice diagnostic on $34{,}858$ communes}

\author[1]{Rohan Foss\'e\thanks{Corresponding author: \texttt{rfosse@cesi.fr}. ORCID: \href{https://orcid.org/0009-0002-2195-0198}{0009-0002-2195-0198}}}
\author[2]{Ga\"el Pallares\thanks{ORCID: \href{https://orcid.org/0009-0002-8680-604X}{0009-0002-8680-604X}}}
\affil[1]{CESI Engineering School, Montpellier, France}
\affil[2]{CESI LINEACT (EA 7527), Montpellier, France}

\date{}

\begin{document}
\maketitle

% =========================================================================
\begin{abstract}
\noindent
The standard French cycling-policy diagnostic answers the question
\emph{where} the cycling environment is weak.  It does not answer
the policy question that follows -- \emph{who pays the cost of that
weakness}.  We close this gap by stacking three vulnerability
layers on top of the commune-level IMD-4 cycling-environment
indicator \citep{Fosse2026imd_national} : cycling-environment
deprivation (bottom-$33\,\%$ IMD-4), monetary-poverty exposure
(top-$33\,\%$ INSEE poverty rate, $> 15\,\%$), and a structural
geographic-isolation layer that captures the joint contribution of
the lowest income tertile and the overseas \emph{d\'epartements}.
The triple-penalty intersection identifies \textbf{322 of the
$362$ ``double-penalty'' cycling-poverty deserts} as also
monetary-poverty-vulnerable, covering $1.89$ million inhabitants
or $91.7\,\%$ of the desert population.  $42$ of the $362$ deserts
are in the overseas \emph{d\'epartements} -- $11.6\,\%$ of the
count but $36.3\,\%$ of the population -- making
Outre-mer a structurally over-represented mobility-justice
priority.  A subset of \textbf{$90$ metropolitan deserts}, located
in the national first income decile but outside Outre-mer, forms
the policy-blind spot of the current Plan V\'elo distribution :
visible neither under a poverty-rate criterion alone nor under an
Outre-mer flag alone, they aggregate to $\sim 450{,}000$
inhabitants in mid-size urban and peri-urban France.  We argue
that the IMD-4 read jointly with these three vulnerability layers
turns a measurement instrument into a mobility-justice diagnostic
\citep{Pereira2017,Verlinghieri2020}, with a concrete and
actionable list of priority communes for the second tranche of
the Plan V\'elo (2023--2027).

\noindent\textbf{Keywords:}
cycling poverty ; mobility justice ; transport equity ; composite
indicator ; commune-level diagnostic ; French Plan V\'elo ;
overseas territories.
\end{abstract}

% =========================================================================
\section{Introduction}
\label{sec:intro}

The French cycling-policy debate has matured along a measurement
axis : the FUB \emph{Barom\`etre V\'elo} \citep{FUB2023} captures
perception, the EMP modal-share survey \citep{EMP2019} captures
behaviour, the INSEE recensement \emph{part-v\'elo-travail}
\citep{INSEEMobPro2022} captures exhaustive commune-level
outcomes, and a recent composite indicator
\citep{Fosse2026imd_national} integrates these into a
$34{,}858$-commune supply-side score (the IMD-4).  Each step
sharpens the answer to \emph{where} the cycling environment is
weak, and each step contributes to the Plan V\'elo's
\citep{PlanVelo2018} investment-targeting toolkit.

What this measurement-first agenda does not answer is the
distributive question that follows : among the communes that are
cycling-environment-deprived, \emph{which residents pay the
cost of that deprivation}.  Transport-justice scholarship
\citep{Lucas2012,Martens2012,Pereira2017,Sheller2018,Verlinghieri2020}
has long argued that the social geography of mobility deprivation
matters at least as much as its physical geography ; in France,
that joint diagnostic has so far been published only at the
city-level \citep{Cerema2022,Smith2015,Hosford2018} or in
narrative form rather than at commune resolution on a national
panel.

This paper closes the gap by stacking three commune-level
vulnerability layers on top of the IMD-4 and computing the
triple-penalty intersection across all $34{,}858$ French
communes.  We answer three operational questions :

\begin{enumerate}[leftmargin=*]
    \item[\textbf{Q1}] Among the $362$ ``double-penalty'' cycling-
          poverty deserts identified by \citet{Fosse2026imd_national}
          on the joint bottom-IMD / high-poverty criterion, how
          many remain priority once a third vulnerability layer
          (monetary poverty by income tertile, or structural
          Outre-mer geography) is added ?

    \item[\textbf{Q2}] How does the population concentration of
          the desert list shift between the metropolitan and
          overseas \emph{d\'epartements} once the triple-penalty
          intersection is computed ?

    \item[\textbf{Q3}] Which communes are \emph{invisible} to the
          current desert-list framing -- below the cycling
          environment threshold and in the first national income
          decile but neither in Outre-mer nor classified by their
          local poverty rate as ``poor'' -- and how large is this
          policy-blind subset ?
\end{enumerate}

\noindent\textbf{Contributions.}
\textbf{(C1)} A formal triple-penalty index ($\TPI$,
Section~\ref{sec:method}) that overlays three vulnerability
layers (cycling-environment, monetary poverty, structural
geography) on the IMD-4 substrate.
\textbf{(C2)} A commune-level diagnostic that recovers $322$
triple-penalty communes out of the $362$ original deserts
($88.9\,\%$, covering $91.7\,\%$ of desert population), confirms
the structural over-representation of Outre-mer ($11.6\,\%$ of
the count but $36.3\,\%$ of the population), and surfaces a
$90$-commune metropolitan-extreme subset
(Section~\ref{sec:results}).
\textbf{(C3)} An operational reading of the $\TPI$ for the
$2023$--$2027$ tranche of the Plan V\'elo : a ranked, justice-
weighted list of priority communes that complements the
measurement-only ranking of \citet{Fosse2026imd_national}.

% =========================================================================
\section{Data}
\label{sec:data}

\subsection{Cycling-environment substrate}

We take as given the IMD-4 (\emph{Indice de Mobilit\'e Douce \`a
quatre composantes}) computed by \citet{Fosse2026imd_national} on
all $34{,}858$ French communes.  The IMD-4 combines four
supply-side axes -- transit multimodality (M), cycling
infrastructure density (I), topography (T) and population density
(D) -- under a Bayesian simplex prior calibrated on a $59$-city
dock-based bike-sharing panel.  The companion paper establishes
that the IMD-4 (i) beats the de facto national standard, the
Cerema cycling-infrastructure density \citep{Cerema2021Velo}, by
$+18\,$pt $R^2$ in predicting INSEE part-v\'elo-travail, and
(ii) is essentially orthogonal to commune income within the urban
subset ($\rho = +0.001$, $n = 31{,}266$).  We use the
national-application z-scored IMD-4 score and its per-commune
$M^z$, $I^z$, $D^z$ components as input.

\subsection{Social vulnerability layers}

Three commune-level social vulnerability variables are layered on
the IMD-4 :

\begin{description}
    \item[\textbf{Income median}] -- INSEE Filosofi commune-level
          median income per consumption unit (\emph{revenu m\'edian
          disponible par UC}), most recent vintage, in euros per
          year.  Available on $31{,}266$ communes.

    \item[\textbf{Poverty rate}] -- INSEE Filosofi commune-level
          share of the population under the national $60\,\%$-of-
          median poverty threshold.  Available on $4{,}408$ communes
          (the publication-disclosure threshold filters out small
          rural communes).

    \item[\textbf{Outre-mer status}] -- a binary flag indicating
          one of the five overseas \emph{d\'epartements}
          (Guadeloupe, Martinique, Guyane, La R\'eunion, Mayotte,
          INSEE department code starting with $97$).  We treat
          Outre-mer status as a structural inequity layer in line
          with the mobility-justice literature on post-colonial
          geographies \citep{Sheller2018}.
\end{description}

Income median, poverty rate, and IMD-4 are extracted from the
companion-paper public release ; the Outre-mer flag is derived
from the INSEE commune code.

% =========================================================================
\section{Method : the triple-penalty index}
\label{sec:method}

\subsection{Definition}

For each commune $c$ we define three binary deprivation flags :
\begin{align*}
    L_1(c) &= \mathbb{1}\big[\IMD\text{-}4_c < q_{33}(\IMD\text{-}4)\big]
              \quad &&\text{(cycling-environment deprivation)} \\
    L_2(c) &= \mathbb{1}\big[\text{poverty\_rate}_c > 15\,\%\big]
              \quad &&\text{(monetary-poverty exposure)} \\
    L_3(c) &= \mathbb{1}\big[\text{income}_c < q_{33}(\text{income})\big]
              \;\lor\; \mathbb{1}\big[\text{overseas}_c\big]
              \quad &&\text{(structural deprivation)}
\end{align*}
The triple-penalty index is the joint indicator
\[
    \TPI(c) \;=\; L_1(c) \cdot L_2(c) \cdot L_3(c),
\]
and a commune is in the triple-penalty subset whenever
$\TPI(c) = 1$.  $L_1 \land L_2 = 1$ recovers the
$362$-commune ``double-penalty desert'' list of the companion
paper ; the $L_3$ multiplier adds the third layer.

\subsection{Numerical thresholds}

The IMD-4 $33\,\%$-tile threshold $q_{33}(\IMD\text{-}4) = -0.027$
is inherited unchanged.  The poverty threshold $15\,\%$ is the
INSEE Filosofi national reference cut.  The income $33\,\%$-tile
threshold $q_{33}(\text{income}) = 21{,}790$ EUR\,/\,year is
computed on the $31{,}266$ communes with reported median income ;
the first-decile threshold $q_{10}(\text{income}) = 19{,}940$
EUR\,/\,year is used in the metropolitan-extreme subset of
Section~\ref{sec:metro-extreme}.

% =========================================================================
\section{Results}
\label{sec:results}

\subsection{Triple-penalty intersection}

Applied to the $4{,}408$ communes with both IMD-4 and poverty
rate, the triple-penalty intersection identifies
\textbf{$322$ communes}, against the $362$ of the IMD-only desert
list (Table~\ref{tab:tpi}).  $88.9\,\%$ of the original deserts
remain in the triple-penalty subset ; the residual $40$ communes
are those classified as ``poor'' under the local poverty-rate
criterion but whose income median sits above the national lowest
tertile.

\begin{table}[H]
\centering
\caption{Triple-penalty subset relative to the $362$-commune
double-penalty desert list.  Population columns use INSEE
population.}
\label{tab:tpi}
\small
\begin{tabular}{@{}lrrr@{}}
\toprule
Subset & Communes & \% of $362$ & Population \\
\midrule
$L_1 \land L_2$ (cycling-poverty deserts)  & 362 & $100\,\%$  & $2{,}059{,}864$ \\
$\TPI = L_1 \land L_2 \land L_3$           & 322 & $88.9\,\%$ & $1{,}888{,}163$ \\
\quad of which Outre-mer                    &  42 & $11.6\,\%$ &   $748{,}287$ \\
\quad of which metropolitan first decile    &  90 & $24.9\,\%$ &   $\sim 450{,}000$ \\
\bottomrule
\end{tabular}
\end{table}

\subsection{Outre-mer over-representation}

Of the $322$ triple-penalty communes, $42$ are in the overseas
\emph{d\'epartements} : $11.6\,\%$ of the count but
\textbf{$36.3\,\%$ of the population} ($748{,}287$ inhabitants).
The discrepancy reflects the structural demographic profile of
the overseas networks : the Outre-mer towns that fail the IMD-4
threshold are mid-to-large municipalities, against the
metropolitan deserts which are dominated by small peri-urban and
rural communes.  La R\'eunion alone accounts for the top six
deserts by population (Saint-Paul, Le Tampon, Saint-Louis,
Saint-Joseph, Saint-Beno\^it, Sainte-Marie).  The cycling-
environment gap in Outre-mer therefore translates, in absolute
population terms, into a substantially larger affected cohort
than its commune count suggests.

\subsection{Metropolitan-extreme subset : the policy-blind 90}
\label{sec:metro-extreme}

The triple-penalty intersection includes a subset of
\textbf{$90$ metropolitan communes} located in the national
first income decile ($< 19{,}940$\,EUR per UC) but outside the
overseas \emph{d\'epartements}.  These communes are simultaneously
cycling-environment-deprived, locally poverty-exposed, and
nationally income-deprived.  They cluster in two geographies :
mid-size southern cities and their peripheries (Arles,
Saint-Gilles, Berre-l'\'Etang, Graulhet, Carcassonne) where
infrastructure is sparse and transit is thin, and post-industrial
medium cities of the north and east (Autun, Cholet, etc.) where
cycling infrastructure is present but multimodal access has
declined.

The policy reading is that these $90$ communes are
\textbf{simultaneously invisible} under two of the standard
single-axis equity filters : (i) a poverty-rate filter alone, used
by some Plan V\'elo sub-actions, retains them but submerges them
inside the $362$-strong general list ; (ii) an Outre-mer-priority
filter, used by ANCT and DGOM as a coarse structural criterion,
omits them entirely.  Only the joint triple-penalty intersection
surfaces them as a distinct priority cohort, with a cumulative
population of approximately $450{,}000$ inhabitants and a per-
commune median below the desert average.

\subsection{Top of the triple-penalty ranking}

Table~\ref{tab:tpi-top} reports the ten most populous communes
in the triple-penalty subset.  Seven of the ten are in La R\'eunion,
with $\sim 392{,}000$ cumulative inhabitants ; Arles (Bouches-du-
Rh\^one) is the most populous metropolitan triple-penalty case.

\begin{table}[H]
\centering
\caption{Top ten triple-penalty communes by population.
Pop = INSEE municipal population.  Income = INSEE Filosofi
median revenu disponible par UC (EUR/year).  Pov = INSEE Filosofi
poverty rate (\%).  All ten communes satisfy $\TPI = 1$ ; the
seven Outre-mer entries also satisfy the ultra-penalty
intersection $L_1 \land L_2 \land L_3 \land \text{Outre-mer}$.}
\label{tab:tpi-top}
\small
\begin{tabular}{@{}llrrrl@{}}
\toprule
Commune & Code & Pop & Income (EUR/yr) & Pov (\%) & D\'epartement \\
\midrule
Saint-Paul          & 97415 & 108{,}088 & 18{,}910 & 30 & La R\'eunion \\
Le Tampon           & 97422 &  82{,}579 & 16{,}560 & 38 & La R\'eunion \\
Saint-Louis         & 97414 &  54{,}941 & 15{,}270 & 43 & La R\'eunion \\
Arles               & 13004 &  51{,}811 & 20{,}060 & 24 & Bouches-du-Rh\^one \\
Saint-Joseph        & 97412 &  39{,}207 & 15{,}200 & 43 & La R\'eunion \\
Saint-Beno\^it      & 97410 &  38{,}604 & 14{,}550 & 46 & La R\'eunion \\
Sainte-Marie        & 97418 &  37{,}093 & 19{,}110 & 30 & La R\'eunion \\
La Possession       & 97408 &  36{,}568 & 20{,}980 & 24 & La R\'eunion \\
Saint-Leu           & 97413 &  36{,}163 & 17{,}630 & 34 & La R\'eunion \\
Sainte-Suzanne      & 97420 &  25{,}551 & 17{,}720 & 34 & La R\'eunion \\
\bottomrule
\end{tabular}
\end{table}

% =========================================================================
\section{Discussion}
\label{sec:discussion}

\subsection{From measurement to mobility-justice diagnostic}

The IMD-4 of \citet{Fosse2026imd_national} provides a defensible
answer to the question \emph{where} the French cycling environment
is weak.  Layered with the three vulnerability filters of
Section~\ref{sec:method}, it provides an equally defensible answer
to the question \emph{who pays the cost} of that weakness.  The
triple-penalty intersection is not a re-weighting of the IMD-4 ;
it is a complementary reading that converts the indicator from a
measurement instrument into an actionable mobility-justice
diagnostic in the sense of
\citet{Pereira2017} and \citet{Verlinghieri2020}.  The headline
finding -- $\sim 92\,\%$ of the desert-list population is also
monetary-poverty-vulnerable -- argues that the cycling-environment
gap and the broader transport-poverty literature describe the
same population.

\subsection{Outre-mer is structurally over-represented, not
incidentally}

The Outre-mer concentration of triple-penalty communes
($11.6\,\%$ of the count, $36.3\,\%$ of the population) is not a
quirk of the indicator but a structural feature of the
mobility-supply geography of the overseas \emph{d\'epartements}.
La R\'eunion is mountainous, mid-density, and served by a single
non-electrified rail line ; the M axis of the IMD-4 is therefore
near-zero across most of the island, and the I axis reflects a
recent and incomplete cycling-infrastructure rollout.  These
exposures are well-documented in the mobility-justice literature
\citep{Sheller2018}.  The contribution of this paper is to make
them \emph{commune-level operational} : the triple-penalty list
provides a per-commune target rather than a region-level
narrative.

\subsection{The metropolitan-extreme blind spot}

The $90$ metropolitan-extreme communes
(Section~\ref{sec:metro-extreme}) are the most consequential
policy finding of this paper.  They are simultaneously absent
from the two standard Plan V\'elo equity filters (poverty-rate
alone, Outre-mer flag alone) yet they satisfy a strict
intersection of three deprivation criteria.  If the second tranche
of the Plan V\'elo ($2023$--$2027$) prioritises by either of the
single-axis filters alone, this cohort is materially under-served
even though its triple-penalty profile is among the strongest in
the national distribution.  The $\TPI$ framing is the simplest
modification that surfaces them.

\subsection{Limitations and follow-up work}

\textbf{Missing vulnerability layers.}  The third $L_3$ layer
combines income and Outre-mer status as a structural deprivation
proxy.  A richer specification would add the share of car-less
households, the share of single-parent families, the share of
residents aged $\geq 65$, and the female-employment-rate gap,
all of which the cycling-equity literature documents as bias
amplifiers \citep{Garrard2008,Aldred2016,Mattioli2017}.  These
variables are available in INSEE Filosofi at commune level and
would extend the $L_3$ specification ; we leave their integration
to a follow-up release.

\textbf{Static.}  The triple-penalty intersection is computed on
a single cross-section.  A longitudinal version tracking the
$\TPI$ between successive INSEE waves would let us test whether
Plan V\'elo investment in $\TPI$-flagged communes correlates with
subsequent cycling-environment improvement
\citep{Abadie2010,Abadie2021}.

\textbf{Outcome-side blindness.}  The diagnostic is supply-side
plus social-exposure ; it does not capture the cycling-uptake
behaviour of the affected residents.  Triangulating the $\TPI$
with INSEE part-v\'elo-travail \citep{INSEEMobPro2022} would
quantify the gap between the cycling-environment provided by the
state and the cycling-behaviour observed, by triple-penalty
status.

% =========================================================================
\section{Conclusion}
\label{sec:conclusion}

We introduced the triple-penalty index ($\TPI$), a commune-level
mobility-justice diagnostic that overlays cycling-environment
deprivation, local monetary-poverty exposure, and structural
geographic vulnerability on the IMD-4 substrate of
\citet{Fosse2026imd_national}.  Applied to the $34{,}858$ French
communes, the $\TPI$ recovers $322$ of the $362$ original
``double-penalty'' deserts, confirms the structural over-
representation of Outre-mer ($11.6\,\%$ of the count and
$36.3\,\%$ of the population), and surfaces a $90$-commune
metropolitan-extreme subset that is simultaneously invisible
under the standard single-axis equity filters.

The contribution is not a new indicator but a new reading of
existing public data.  The cycling-environment indicator and the
INSEE Filosofi social variables are already on the open-data
portals ; what was missing was their joint, commune-resolution
intersection.  The triple-penalty list is the first operational
mobility-justice product for the French cycling-policy debate at
the granularity at which local investment decisions are actually
taken.

\bibliography{references_social}

\end{document}
